% Created 2022-07-20 mié 20:16
% Intended LaTeX compiler: pdflatex
\documentclass[11pt]{article}
\usepackage[utf8]{inputenc}
\usepackage[T1]{fontenc}
\usepackage{graphicx}
\usepackage{grffile}
\usepackage{longtable}
\usepackage{wrapfig}
\usepackage{rotating}
\usepackage[normalem]{ulem}
\usepackage{amsmath}
\usepackage{textcomp}
\usepackage{amssymb}
\usepackage{capt-of}
\usepackage{hyperref}
\usepackage[spanish]{babel}
\documentclass[10pt,a5paper]{article}
\usepackage[margin=1cm,a5paper]{geometry}
\pagenumbering{gobble}
\setcounter{secnumdepth}{1}
\author{Luis Eduardo Galindo Amaya}
\date{\today}
\title{Ejemplo de Quine–McCluskey}
\hypersetup{
 pdfauthor={Luis Eduardo Galindo Amaya},
 pdftitle={Ejemplo de Quine–McCluskey},
 pdfkeywords={},
 pdfsubject={},
 pdfcreator={Emacs 26.3 (Org mode 9.1.9)}, 
 pdflang={Spanish}}
\begin{document}

\maketitle
\tableofcontents

\toc

\newpage

\section{Control De Un Display De Siete Segmentos}
\label{sec:orgbf0d6db}
Cada variable representa la activacion de cada uno de los segmentos del 
display, 'BCD' es el valor codificado en binario del numero, entonces
lo que estamos haciendo es convertir cuetro entradas binarias en siete
para controlar el display.

\vspace{0.5cm}

\noindent
\begin{minipage}{.5\textwidth}


\begin{center}
\begin{tabular}{|c|c|ccccccc|}
\hline
ID & BCD & A & B & C & D & E & F & G\\
\hline
0 & 0000 & 1 & 1 & 1 & 1 & 1 & 1 & 0\\
1 & 0001 & 0 & 1 & 1 & 0 & 0 & 0 & 0\\
2 & 0010 & 1 & 1 & 0 & 1 & 1 & 0 & 1\\
3 & 0011 & 1 & 1 & 1 & 1 & 0 & 0 & 1\\
4 & 0100 & 0 & 1 & 1 & 0 & 0 & 1 & 0\\
5 & 0101 & 1 & 0 & 1 & 1 & 0 & 1 & 1\\
6 & 0110 & 1 & 0 & 1 & 1 & 1 & 1 & 1\\
7 & 0111 & 1 & 1 & 1 & 0 & 0 & 0 & 0\\
8 & 1000 & 1 & 1 & 1 & 1 & 1 & 1 & 1\\
9 & 1001 & 1 & 1 & 1 & 0 & 0 & 1 & 1\\
10 & 1010 & x & x & x & x & x & x & x\\
11 & 1011 & x & x & x & x & x & x & x\\
12 & 1100 & x & x & x & x & x & x & x\\
13 & 1101 & x & x & x & x & x & x & x\\
14 & 1110 & x & x & x & x & x & x & x\\
15 & 1111 & x & x & x & x & x & x & x\\
\hline
\end{tabular}
\end{center}
\end{minipage}
\begin{minipage}{.5\textwidth}


\begin{center}
\includegraphics[width=3.5cm]{./7_Segment_Display.png}
\end{center}
\end{minipage}

\section{Variable A}
\label{sec:org554a0db}
Cada columna de la tabla corresponde a una variable, sí extraemos la columna
correspondiente a 'A' obtendremos algo como esto:

\begin{center}
\begin{tabular}{|ccccccccccccccccc|}
\hline
ID & 0 & 1 & 2 & 3 & 4 & 5 & 6 & 7 & 8 & 9 & 10 & 11 & 12 & 13 & 14 & 15\\
\hline
 & 1 & 0 & 1 & 1 & 0 & 1 & 1 & 1 & 1 & 1 & x & x & x & x & x & x\\
\hline
\end{tabular}
\end{center}

Podemos convertir esta tabla en su notacion de sumatoria equivalente para 
ingresarla mas facilmente en el programa, de igual manera el programa tiene
soporte para ingresar tablas:

\begin{equation}
A(w,x,y,z) = \(\sum\) m(0,2,3,5,6,7,8,9) + \(\sum\) d(10,11,12,13,14,15)
\label{orga6e9a70}
\end{equation}

los valores en 'm' son las activaciones y 'd' son los valores que no importan,
también conocidos como 'redundancias'. Ahora podemos emepzar a usar el programa.

\subsection*{Entradas}
\label{sec:orgb9a4b97}
\begin{verbatim}
---------------------
      REDUCCION      
  DE FUNCION LOGICA  
---------------------

1.LISTA // 2.TABLA?      Aquí podemos elegir el metodo de entrada
1  
NUMERO DE VARIABLES?     Tenemos 4 variables 'w', 'x', 'y' y 'z'
4
ACTIVACION?              Las activaciones, esto es 'm'
0,2,3,5,6,7,8,9 
REDUNDANCIAS?            Redundancias, esto es 'd'
10,11,12,13,14,15

[0, 2, 3, 5, 6, 7, 8, 9] Aquí muestra la lista 'm'
[10, 11, 12, 13, 14, 15] y aca la lista 'd'

---------------------
      RESULTADO      
---------------------

ESENCIALES:               Los minterminos esenciales
-0-0 (0, 8, 2, 10)                 
--1- (2, 3, 6, 7, 10, 11, 14, 15)
-1-1 (13, 15, 5, 7)
1--- (8, 9, 10, 11, 12, 13, 14, 15)

NO ESENCIALES: []

\end{verbatim}

\newpage

\subsection*{Interpretando los resultados}
\label{sec:org02e39f3}
Ahora tenemos que extraer una expresión booleana de la salida del programa,
debemos tener en cuenta que nuestras variables tienen mas valor dependiendo 
su posicion. En la forma de \hyperref[orga6e9a70]{sumatoria} 'w' es el bit más significativo porque 
esta primero y 'z' es el bit menos significativo, representa el uno. Cada bit
tiene un valor que representa el estado de la variable ('1', '0' ó '-'), '1' 
es activo, '0' es negado  y '-' es que no forma parte del mintermino.

\begin{verbatim}

ESENCIALES:              Los minterminos esenciales
-0-0 (0, 8, 2, 10)                 
--1- (2, 3, 6, 7, 10, 11, 14, 15)
-1-1 (13, 15, 5, 7)
1--- (8, 9, 10, 11, 12, 13, 14, 15)

NO ESENCIALES: []        Si no hay términos faltantes significa
                         que solo existe una solución para nuestras
                         especificaciones

\end{verbatim}

Ahora como todos los mintermino son esenciales, lo cual significa que es 
imposible omitirlos para obtener nuestra funcion, los unimos para crear la 
función y obtenemos:

\begin{center}
\begin{tabular}{llllll}
w & x & y & z &  & Representación\\
\hline
- & 0 & - & 0 &  & \(\bar x \bar z\)\\
- & - & 1 & - &  & \(y\)\\
- & 1 & - & 1 &  & \(xz\)\\
1 & - & - & - &  & \(w\)\\
\hline
\end{tabular}
\end{center}

\begin{equation}
A(w,x,y,z) = xz + \bar x \bar z + y + w
\end{equation}

Esta es la forma de suma mas simplificada de la funcion, si usamos un poco de 
álgebra booleana podemos reducir un poco más la funcion añadiendo una compuerta
'xnor':

\begin{equation}
A(w,x,y,z) = x \(\odot\) z + y + w
\end{equation}
\section{Variable B}
\label{sec:orgd0bb2bd}
\subsection*{Activaciones}
\label{sec:orgde64057}
\begin{center}
\begin{tabular}{|ccccccccccccccccc|}
\hline
ID & 0 & 1 & 2 & 3 & 4 & 5 & 6 & 7 & 8 & 9 & 10 & 11 & 12 & 13 & 14 & 15\\
\hline
 & 1 & 1 & 1 & 1 & 1 & 0 & 0 & 1 & 1 & 1 & x & x & x & x & x & x\\
\hline
\end{tabular}
\end{center}
\begin{equation}
B(w,x,y,z) = \(\sum\) m(0,1,2,3,4,7,8,9) + \(\sum\) d(10,11,12,13,14,15)
\end{equation}

\subsection*{Entradas}
\label{sec:orgae988ba}
\begin{verbatim}
---------------------
      REDUCCION      
  DE FUNCION LOGICA  
---------------------

1.LISTA // 2.TABLA?
1
NUMERO DE VARIABLES?
4
ACTIVACION?
0,1,2,3,4,7,8,9
REDUNDANCIAS?
10,11,12,13,14,15

---------------------
      RESULTADO      
---------------------

ESENCIALES:
-0-- (0, 1, 2, 3, 8, 9, 10, 11)
--00 (0, 8, 4, 12)
--11 (11, 3, 15, 7)

NO ESENCIALES: []                    
1--- (8, 9, 10, 11, 12, 13, 14, 15)  Hay términos no esenciales 
                                     pero no hay términos 
                                     faltantes así que se ignora
\end{verbatim}

\section{Variable C}
\label{sec:org7afeeb1}
\subsection*{Activaciones}
\label{sec:org03d52f7}
\begin{center}
\begin{tabular}{|ccccccccccccccccc|}
\hline
ID & 0 & 1 & 2 & 3 & 4 & 5 & 6 & 7 & 8 & 9 & 10 & 11 & 12 & 13 & 14 & 15\\
\hline
 & 1 & 1 & 0 & 1 & 1 & 1 & 1 & 1 & 1 & 1 & x & x & x & x & x & x\\
\hline
\end{tabular}
\end{center}
\begin{equation}
C(w,x,y,z) = \(\sum\) m(0,1,3,4,5,6,7,8,9) + \(\sum\) d(10,11,12,13,14,15)
\end{equation}

\subsection*{Entradas}
\label{sec:org4b863c4}
\begin{verbatim}
---------------------
      REDUCCION      
  DE FUNCION LOGICA  
---------------------

1.LISTA // 2.TABLA?
1
NUMERO DE VARIABLES?
4
ACTIVACION?
0,1,3,4,5,6,7,8,9
REDUNDANCIAS?
10,11,12,13,14,15

[0, 1, 3, 4, 5, 6, 7, 8, 9]
[10, 11, 12, 13, 14, 15]

---------------------
      RESULTADO      
---------------------

ESENCIALES:
--0- (0, 1, 4, 5, 8, 9, 12, 13)
---1 (1, 3, 5, 7, 9, 11, 13, 15)
-1-- (4, 5, 6, 7, 12, 13, 14, 15)

NO ESENCIALES: []
1--- (8, 9, 10, 11, 12, 13, 14, 15)
\end{verbatim}
\section{Variable D}
\label{sec:org1357122}
\subsection*{Activaciones}
\label{sec:org29f051e}
\begin{center}
\begin{tabular}{|ccccccccccccccccc|}
\hline
ID & 0 & 1 & 2 & 3 & 4 & 5 & 6 & 7 & 8 & 9 & 10 & 11 & 12 & 13 & 14 & 15\\
\hline
 & 1 & 0 & 1 & 1 & 0 & 1 & 1 & 0 & 1 & 0 & x & x & x & x & x & x\\
\hline
\end{tabular}
\end{center}
\begin{equation}
D(w,x,y,z) = \(\sum\) m(0,2,3,5,6,8) + \(\sum\) d(10,11,12,13,14,15)
\end{equation}

\subsection*{Entradas}
\label{sec:org11e9501}
\begin{verbatim}
---------------------
      REDUCCION      
  DE FUNCION LOGICA  
---------------------

1.LISTA // 2.TABLA?
1
NUMERO DE VARIABLES?
4
ACTIVACION?
0,2,3,5,6,8
REDUNDANCIAS?
10,11,12,13,14,15

---------------------
      RESULTADO      
---------------------

ESENCIALES:
-0-0 (0, 8, 2, 10)
-01- (10, 11, 2, 3)
-101 (13, 5)
--10 (2, 10, 6, 14)

NO ESENCIALES: []
1--0 (8, 10, 12, 14)
1-1- (10, 11, 14, 15)
11-- (12, 13, 14, 15)
\end{verbatim}

\section{Variable E}
\label{sec:org1b4d3c0}
\subsection*{Activaciones}
\label{sec:org8e9078b}
\begin{center}
\begin{tabular}{|ccccccccccccccccc|}
\hline
ID & 0 & 1 & 2 & 3 & 4 & 5 & 6 & 7 & 8 & 9 & 10 & 11 & 12 & 13 & 14 & 15\\
\hline
 & 1 & 0 & 1 & 0 & 0 & 0 & 1 & 0 & 1 & 0 & x & x & x & x & x & x\\
\hline
\end{tabular}
\end{center}
\begin{equation}
E(w,x,y,z) = \(\sum\) m(0,2,6,8) + \(\sum\) d(10,11,12,13,14,15)
\end{equation}

\subsection*{Entradas}
\label{sec:org24b0af6}
\begin{verbatim}
---------------------
      REDUCCION      
  DE FUNCION LOGICA  
---------------------

1.LISTA // 2.TABLA?
1
NUMERO DE VARIABLES?
4      
ACTIVACION?
0,2,6,8
REDUNDANCIAS?
10,11,12,13,14,15

[0, 8, 2, 6]
[10, 11, 12, 13, 14, 15]

---------------------
      RESULTADO      
---------------------

ESENCIALES:
-0-0 (0, 8, 2, 10)
--10 (2, 10, 6, 14)

NO ESENCIALES: []
1--0 (8, 10, 12, 14)
1-1- (10, 11, 14, 15)
11-- (12, 13, 14, 15)
\end{verbatim}

\section{Variable F}
\label{sec:orgd672bc5}
\subsection*{Activaciones}
\label{sec:org23fe14e}
\begin{center}
\begin{tabular}{|ccccccccccccccccc|}
\hline
ID & 0 & 1 & 2 & 3 & 4 & 5 & 6 & 7 & 8 & 9 & 10 & 11 & 12 & 13 & 14 & 15\\
\hline
 & 1 & 0 & 0 & 0 & 1 & 1 & 1 & 0 & 1 & 1 & x & x & x & x & x & x\\
\hline
\end{tabular}
\end{center}
\begin{equation}
F(w,x,y,z) = \(\sum\) m(0,4,5,6,6,9) + \(\sum\) d(10,11,12,13,14,15)
\end{equation}

\subsection*{Entradas}
\label{sec:org20eb8ad}
\begin{verbatim}
---------------------
      REDUCCION      
  DE FUNCION LOGICA  
---------------------

1.LISTA // 2.TABLA?
1
NUMERO DE VARIABLES?
4
ACTIVACION?
0,4,5,6,6,9
REDUNDANCIAS?
10,11,12,13,14,15

[0, 4, 5, 6, 9]
[10, 11, 12, 13, 14, 15]

---------------------
      RESULTADO      
---------------------

ESENCIALES:
0-00 (0, 4)
-10- (12, 13, 4, 5)
-1-0 (4, 12, 6, 14)
1--1 (9, 11, 13, 15)

NO ESENCIALES: []
1-1- (10, 11, 14, 15)
11-- (12, 13, 14, 15)
\end{verbatim}

\section{Variable G}
\label{sec:org7b783ff}
\subsection*{Activaciones}
\label{sec:org01bde96}
\begin{center}
\begin{tabular}{|ccccccccccccccccc|}
\hline
ID & 0 & 1 & 2 & 3 & 4 & 5 & 6 & 7 & 8 & 9 & 10 & 11 & 12 & 13 & 14 & 15\\
\hline
 & 0 & 0 & 1 & 1 & 0 & 1 & 1 & 0 & 1 & 1 & x & x & x & x & x & x\\
\hline
\end{tabular}
\end{center}
\begin{equation}
G(w,x,y,z) = \(\sum\) m(2,3,5,6,8,9) + \(\sum\) d(10,11,12,13,14,15)
\end{equation}

\subsection*{Entradas}
\label{sec:orgd815a60}
\begin{verbatim}
---------------------
      REDUCCION      
  DE FUNCION LOGICA  
---------------------

1.LISTA // 2.TABLA?
1
NUMERO DE VARIABLES?
4
ACTIVACION?
2,3,5,6,8,9
REDUNDANCIAS?
10,11,12,13,14,15

[2, 3, 5, 6, 8, 9]
[10, 11, 12, 13, 14, 15]

---------------------
      RESULTADO      
---------------------

ESENCIALES:
-01- (10, 11, 2, 3)
-101 (13, 5)
--10 (2, 10, 6, 14)
1--- (8, 9, 10, 11, 12, 13, 14, 15)

NO ESENCIALES: []
\end{verbatim}

\newpage

\section{Resultados}
\label{sec:org44d47ec}



\begin{equation}
  \begin{aligned}
    A &: ~ xz + \bar x \bar z + y + w \\
    B &: ~ \bar y \bar z + yz + \bar x \\
    C &: ~ \bar y + z + x \\
    D &: ~ x \bar y z + \bar x y + y \bar z + \bar x  \bar z \\
    E &: ~ \bar x \bar z + y \bar z \\
    F &: ~ \bar w \bar y \bar z + x \bar y + x \bar z + wz \\
    G &: ~ \bar x y + x \bar y z + y \bar z + w
  \end{aligned}
\end{equation}
\end{document}
